\documentclass[11pt]{article}
\usepackage[utf8]{inputenc}
\usepackage{geometry}
\usepackage{hyperref}
\usepackage{enumitem}

\geometry{
    a4paper,
    left=1in,
    right=1in,
    bottom=1in,
    top=0.5in
}

\title{\vspace{-1.5cm}String Art Generator}
\author{Project Documentation}

\begin{document}

\maketitle

\section{Goal}
The String Art Generator project aims to transform digital images into string art representations by simulating the process of wrapping a single continuous string around a set of fixed points (nails) arranged in a circle.

\section{Technical Approach}

\subsection{Core Algorithm}
The project implements a greedy algorithm in the \texttt{StringRing} Rust structure, which:
\begin{itemize}
    \item Converts input images to grayscale matrices
    \item Maintains two key matrices:
        \begin{itemize}
            \item \texttt{input\_image}: Original grayscale values
            \item \texttt{drawn\_image}: Current state of string art
        \end{itemize}
    \item Uses Bresenham's line algorithm for accurate line drawing
    \item Implements a greedy selection strategy:
        \begin{enumerate}
            \item Calculates potential impact of all possible string placements
            \item Selects the string placement that maximizes darkness in areas matching the input image
            \item Updates the drawn image state and pixel weights
        \end{enumerate}
\end{itemize}

\subsection{Implementation}
The system is structured in multiple layers:
\begin{itemize}
    \item \textbf{Core Library} (\texttt{stringart}) that implements StringRing algorithm:
    \item \textbf{WebAssembly Interface} (\texttt{wasm-stringart})
    \item \textbf{Desktop Application} (\texttt{src-tauri})(deprecated)
\end{itemize}

\section{Evaluation}

The project can be evaluated through its web interface, which provides several key metrics:


The web interface (\url{https://omar-abdelgawad.github.io/string-art-app/}) provides an accessible platform for users to evaluate the algorithm's effectiveness and experiment with different parameters to achieve optimal results.

\end{document} 